\section{Two-qubit dephasing and continuous monitoring}
\label{sec:two_qubit_dephasing}

In this section we consider a simple but paradigmatic open-system model: two qubits undergoing pure dephasing along the $z$ axis, and we discuss how different continuous measurement schemes on the environment (photodetection, homodyne and heterodyne detection) give rise to different stochastic unravellings of the same Lindblad dynamics.  
Although the unconditional master equation is the same for all these schemes, the corresponding conditional evolutions and measurement records have very different informational content, with relevant consequences for quantum control, entanglement generation and metrology.

Throughout this section we adopt the convention introduced in the previous chapter: $\ket{0}$ denotes the ground state and $\ket{1}$ the excited state of each qubit, with
\begin{equation}
	\sigma_z \ket{1} = + \ket{1},
	\qquad
	\sigma_z \ket{0} = - \ket{0},
\end{equation}
and we use the ordered computational basis
\begin{equation}
	\bigl\{ \ket{11}, \ket{10}, \ket{01}, \ket{00} \bigr\}.
\end{equation}

\subsection{Unconditional dephasing dynamics for two qubits}
\label{subsec:two_qubit_unconditional}

We consider two qubits labelled by $j=1,2$, each undergoing pure dephasing along the $z$ axis, induced by independent Markovian environments.  
The free Hamiltonian of the qubits can be written as
\begin{equation}
	H_S = \frac{\omega_1}{2}\,\sigma_z^{(1)} + \frac{\omega_2}{2}\,\sigma_z^{(2)} + H_{\mathrm{int}},
\end{equation}
where $\sigma_z^{(j)}$ denotes the Pauli operator acting on qubit $j$, $\omega_j$ are the local transition frequencies and $H_{\mathrm{int}}$ accounts for any coherent interaction between the qubits (for instance an Ising coupling).  
The coupling to the environment is chosen such that no energy is exchanged with the bath: the only effect is the randomization of the relative phases between eigenstates of $\sigma_z^{(j)}$.

At the level of the reduced dynamics, this situation is described by the Lindblad master equation
\begin{equation}
	\frac{\mathrm{d}\hat{\rho}_t}{\mathrm{d}t}
	= - i [H_S, \hat{\rho}_t]
	+ \frac{\kappa}{2}\sum_{j=1}^{2} \mathcal{D}[\sigma_z^{(j)}]\hat{\rho}_t,
	\label{eq:two_qubit_dephasing_master}
\end{equation}
where $\kappa$ is the dephasing rate (assumed equal for the two qubits, for simplicity), and
\begin{equation}
	\mathcal{D}[L]\hat{\rho} = L \hat{\rho} L^{\dagger}
	- \frac{1}{2}\left( L^{\dagger} L \hat{\rho} + \hat{\rho} L^{\dagger} L \right)
	\label{eq:dissipator_definition}
\end{equation}
is the standard Lindblad dissipator.  
Since $\sigma_z^{(j)\,2} = \mathbb{1}$, the dissipator simplifies to
\begin{equation}
	\mathcal{D}[\sigma_z^{(j)}]\hat{\rho}
	= \sigma_z^{(j)} \hat{\rho} \sigma_z^{(j)} - \hat{\rho},
\end{equation}
so that the noise term in Eq.~\eqref{eq:two_qubit_dephasing_master} can be written explicitly as
\begin{equation}
	\frac{\kappa}{2}\sum_{j=1}^{2}
	\mathcal{D}[\sigma_z^{(j)}]\hat{\rho}_t
	=
	\frac{\kappa}{2}\sum_{j=1}^{2} 
	\bigl(
	\sigma_z^{(j)} \hat{\rho}_t \sigma_z^{(j)} - \hat{\rho}_t
	\bigr).
	\label{eq:two_qubit_dephasing_dissipator_simple}
\end{equation}

Equation~\eqref{eq:two_qubit_dephasing_master} preserves the populations in the computational basis and damps the coherences.  
To see this more explicitly, let us expand the density operator as
\begin{equation}
	\hat{\rho}_t = \sum_{\mu,\nu \in \{11,10,01,00\}} 
	\rho_{\mu\nu}(t)\, \ket{\mu}\bra{\nu}.
\end{equation}
Using the fact that each $\sigma_z^{(j)}$ acts diagonally in this basis, it is easy to verify that
\begin{equation}
	\frac{\mathrm{d}}{\mathrm{d}t} \rho_{\mu\mu}(t)
	= - i \bigl( E_{\mu} - E_{\mu} \bigr) \rho_{\mu\mu}(t) = 0,
\end{equation}
where $E_{\mu}$ denotes the energy of $\ket{\mu}$ under $H_S$, while the off-diagonal elements decay as
\begin{equation}
	\frac{\mathrm{d}}{\mathrm{d}t} \rho_{\mu\nu}(t)
	= - i (E_{\mu} - E_{\nu}) \rho_{\mu\nu}(t)
	- \Gamma_{\mu\nu} \rho_{\mu\nu}(t),
	\qquad \mu \neq \nu,
\end{equation}
with a dephasing rate $\Gamma_{\mu\nu}$ that depends on how many local $\sigma_z^{(j)}$ eigenvalues differ between $\ket{\mu}$ e $\ket{\nu}$.  
If the two basis states differ on a single qubit only, the corresponding coherence decays with rate $\kappa$; if they differ on both qubits, the rate is $2\kappa$.  
For instance, the coherence between $\ket{11}$ and $\ket{10}$ is suppressed at rate $\kappa$, whereas the coherence between $\ket{11}$ and $\ket{00}$ decays at rate $2\kappa$.

The unconditional dephasing dynamics hence transforms any initial state toward a classical mixture in the computational basis, without generating entanglement and without modifying the populations.  
However, as we now discuss, different ways of continuously monitoring the dephasing environment lead to very different stochastic descriptions of this same Lindblad evolution.

\subsection{Photodetection unravelling}
\label{subsec:two_qubit_pd}

A conceptually simple microscopic model for pure dephasing consists in coupling each qubit to an independent bosonic mode whose number operator is measured, such that phase kicks on the qubits are associated with the emission of quanta into the environment.  
At the level of the master equation~\eqref{eq:two_qubit_dephasing_master}, this corresponds to collapse operators
\begin{equation}
	L_j = \sqrt{\frac{\kappa}{2}}\, \sigma_z^{(j)},
	\qquad j = 1,2.
\end{equation}
If the output fields are continuously monitored by ideal photon counters with (overall) detection efficiency $\eta \in [0,1]$, the conditional evolution of the system state is governed by a stochastic master equation with Poissonian increments,
\begin{equation}
	\begin{aligned}
		\mathrm{d}\hat{\rho}_t^{(c)}
		&= - i [H_S, \hat{\rho}_t^{(c)}]\, \mathrm{d}t
		+ (1-\eta)\frac{\kappa}{2}\sum_{j=1}^{2} 
		\mathcal{D}[\sigma_z^{(j)}]\hat{\rho}_t^{(c)} \,\mathrm{d}t \\
		&\quad + \sum_{j=1}^{2}
		\left(
		\frac{\sigma_z^{(j)} \hat{\rho}_t^{(c)} \sigma_z^{(j)}}
		{\mathrm{Tr}[\hat{\rho}_t^{(c)}]}
		- \hat{\rho}_t^{(c)}
		\right)
		\bigl(\mathrm{d}N_j(t) - \tfrac{\eta\kappa}{2} \,\mathrm{d}t \bigr),
	\end{aligned}
	\label{eq:two_qubit_pd_sme}
\end{equation}
where $\mathrm{d}N_j(t)$ are independent Poisson increments taking values $0$ or $1$, with mean
\begin{equation}
	\mathbb{E}\bigl[\mathrm{d}N_j(t)\bigr] = \frac{\eta\kappa}{2}\, \mathrm{d}t.
\end{equation}
The last line of Eq.~\eqref{eq:two_qubit_pd_sme} describes stochastic jumps associated with detector clicks: whenever $\mathrm{d}N_j(t)=1$, the state undergoes the unitary transformation
\begin{equation}
	\hat{\rho}_t^{(c)} \longrightarrow
	\frac{\sigma_z^{(j)} \hat{\rho}_t^{(c)} \sigma_z^{(j)}}
	{\mathrm{Tr}[\hat{\rho}_t^{(c)}]}.
\end{equation}
Because $\sigma_z^{(j)}$ is unitary and Hermitian, these jumps are strictly unitary operations: they flip the relative phase between eigenstates of $\sigma_z^{(j)}$ but do not change populations or eigenvalues of $\hat{\rho}_t^{(c)}$.

A crucial feature of this unravelling is that the statistics of the Poisson processes $\mathrm{d}N_j(t)$ is \emph{independent} of the system state.  
Indeed, the average click rate $\eta\kappa/2$ does not depend on $\hat{\rho}_t^{(c)}$, because $L_j^{\dagger}L_j \propto \mathbb{1}$.  
As a consequence, the measurement records $N_j(t)$ contain no information about the actual state of the qubits or about the parameter encoded in $H_S$: the only effect of photodetection in this model is to reveal a sequence of random phase flips.  
Upon averaging over all possible detection records, Eq.~\eqref{eq:two_qubit_pd_sme} reduces to the unconditional master equation~\eqref{eq:two_qubit_dephasing_master}.

From the point of view of state engineering and metrology, this photodetection unravelling is therefore rather uninformative in the case of independent dephasing channels: it reproduces the same loss of coherence as the unconditional dynamics, but the measurement record cannot be exploited to recover information or to correct the noise.  
In later sections, when considering different collapse operators (for instance collective or parity-like combinations of $\sigma_z^{(1)}$ and $\sigma_z^{(2)}$), photodetection will instead play a non-trivial role in generating entanglement.

\subsection{Homodyne unravelling}
\label{subsec:two_qubit_homodyne}

A more informative monitoring strategy consists in measuring a field quadrature of each dephasing channel by continuous homodyne detection.  
In this case the environment is modelled as a bosonic field coupled to the system via the same collapse operators $L_j = \sqrt{\kappa/2}\,\sigma_z^{(j)}$, but the output fields are mixed with strong classical local oscillators and the corresponding photocurrents are continuously recorded.

The conditional state of the two qubits then obeys a diffusive stochastic master equation of the form
\begin{equation}
	\begin{aligned}
		\mathrm{d}\hat{\rho}_t^{(c)}
		&= - i [H_S, \hat{\rho}_t^{(c)}]\, \mathrm{d}t
		+ \frac{\kappa}{2}\sum_{j=1}^{2} 
		\mathcal{D}[\sigma_z^{(j)}]\hat{\rho}_t^{(c)} \,\mathrm{d}t \\
		&\quad + \sqrt{\eta\kappa}\sum_{j=1}^{2}
		\mathcal{H}[\sigma_z^{(j)}]\hat{\rho}_t^{(c)} \,\mathrm{d}W_j(t),
	\end{aligned}
	\label{eq:two_qubit_homodyne_sme}
\end{equation}
where $\mathrm{d}W_j(t)$ are independent Wiener increments with zero mean and variance $\mathbb{E}[\mathrm{d}W_j(t)\mathrm{d}W_{j'}(t)] = \delta_{jj'}\, \mathrm{d}t$, and
\begin{equation}
	\mathcal{H}[L]\hat{\rho} =
	L \hat{\rho} + \hat{\rho} L^{\dagger}
	- \mathrm{Tr}\bigl[(L + L^{\dagger})\hat{\rho}\bigr] \hat{\rho}
	\label{eq:innovation_superoperator}
\end{equation}
is the innovation (or measurement) superoperator.  
For Hermitian operators $L = L^{\dagger}$, Eq.~\eqref{eq:innovation_superoperator} simplifies to
\begin{equation}
	\mathcal{H}[L]\hat{\rho} =
	L \hat{\rho} + \hat{\rho} L
	- 2\,\mathrm{Tr}[L \hat{\rho}] \, \hat{\rho}.
\end{equation}

The associated homodyne photocurrents can be written as
\begin{equation}
	I_j(t)\, \mathrm{d}t =
	\sqrt{\eta\kappa}\,\mathrm{Tr}\!\left[\sigma_z^{(j)} \hat{\rho}_t^{(c)}\right] \mathrm{d}t
	+ \mathrm{d}W_j(t),
	\qquad j = 1,2.
	\label{eq:two_qubit_homodyne_current}
\end{equation}
Contrary to the photodetection case, the average value of $I_j(t)$ depends explicitly on the conditional expectation value of $\sigma_z^{(j)}$.  
The measurement record thus contains information about the instantaneous state of each qubit and about any parameter (such as a frequency) encoded in the Hamiltonian $H_S$.

For fixed $\kappa$ and $\eta$, the unconditional dynamics obtained by averaging Eq.~\eqref{eq:two_qubit_homodyne_sme} over all Wiener trajectories is still given by the dephasing master equation~\eqref{eq:two_qubit_dephasing_master}.  
However, the conditional evolution is qualitatively different:  
the homodyne term drives the state toward eigenstates of $\sigma_z^{(j)}$, resulting in a gradual localization in the computational basis along individual trajectories, rather than in a simple exponential damping of coherences.  
In the two-qubit case, this corresponds to a continuous, stochastic collapse toward the four basis states $\ket{11}$, $\ket{10}$, $\ket{01}$ and $\ket{00}$, with probabilities dictated by the initial state and by the noise.

This “informative” nature of the homodyne record makes it possible, at least in principle, to use continuous measurement and classical post-processing (or feedback) to partially compensate the effect of dephasing noise, or to engineer specific conditional states.  
In particular, when the collapse operators are chosen to be collective combinations of $\sigma_z^{(1)}$ and $\sigma_z^{(2)}$, the homodyne unravelling can be exploited to generate and stabilize entangled states, as will be discussed later.

\subsection{Heterodyne unravelling}
\label{subsec:two_qubit_heterodyne}

Finally, we consider the case in which each dephasing channel is monitored by double-homodyne (heterodyne) detection.  
Operationally, the output field of each channel is split on a balanced beam splitter and each arm is measured by homodyne detection with a relative phase shift of $\pi/2$ between the two local oscillators.  
This configuration enables the simultaneous measurement of two conjugate quadratures of the field, at the cost of adding extra vacuum noise to the measurement record.

For the present model, where the collapse operators are Hermitian and proportional to $\sigma_z^{(j)}$, only one specific quadrature actually carries information about the system observables; the orthogonal quadrature contributes purely vacuum noise.  
As a result, the heterodyne unravelling can be understood as an ``over-resolved'' version of homodyne detection, effectively equivalent to a homodyne measurement with a reduced efficiency.

At the level of the SME, one convenient description of heterodyne detection is
\begin{equation}
	\begin{aligned}
		\mathrm{d}\hat{\rho}_t^{(c)}
		&= - i [H_S, \hat{\rho}_t^{(c)}]\, \mathrm{d}t
		+ \frac{\kappa}{2}\sum_{j=1}^{2} 
		\mathcal{D}[\sigma_z^{(j)}]\hat{\rho}_t^{(c)} \,\mathrm{d}t \\
		&\quad + \sqrt{\frac{\eta\kappa}{2}} \sum_{j=1}^{2}
		\mathcal{H}[\sigma_z^{(j)}]\hat{\rho}_t^{(c)} \,\mathrm{d}W_j^{(1)}(t),
	\end{aligned}
	\label{eq:two_qubit_heterodyne_sme}
\end{equation}
where $\mathrm{d}W_j^{(1)}(t)$ are again Wiener increments.  
Formally this has the same structure as the homodyne SME in Eq.~\eqref{eq:two_qubit_homodyne_sme}, but with an effective measurement strength reduced by a factor $1/\sqrt{2}$, reflecting the fact that only one of the two measured quadratures carries information about $\sigma_z^{(j)}$, while the other contributes additional vacuum fluctuations.

Correspondingly, the two measured currents for each channel can be written as
\begin{equation}
	\begin{aligned}
		I_j^{(1)}(t)\, \mathrm{d}t
		&= \sqrt{\frac{\eta\kappa}{2}}\, 
		\mathrm{Tr}\!\left[\sigma_z^{(j)} \hat{\rho}_t^{(c)}\right] \mathrm{d}t
		+ \mathrm{d}W_j^{(1)}(t), \\
		I_j^{(2)}(t)\, \mathrm{d}t
		&= \text{(pure vacuum noise)},
	\end{aligned}
	\label{eq:two_qubit_heterodyne_currents}
\end{equation}
where $I_j^{(2)}(t)$ does not contain any information about the system state in the present model.  

As in the previous cases, averaging Eq.~\eqref{eq:two_qubit_heterodyne_sme} over all stochastic trajectories recovers the unconditional dephasing dynamics~\eqref{eq:two_qubit_dephasing_master}.  
However, for a fixed physical dephasing rate $\kappa$ and detection efficiency $\eta$, the heterodyne measurement extracts information about the system more slowly than the homodyne one, due to the effective reduction of the signal-to-noise ratio in the informative current $I_j^{(1)}(t)$.  
Conditional states generated by heterodyne monitoring are therefore, on average, less pure than those produced by homodyne monitoring, and the corresponding measurement records are less useful for estimation and control purposes.

In summary, photodetection, homodyne and heterodyne detection of the same two-qubit dephasing channels all induce the same Lindblad master equation at the unconditional level, but correspond to very different stochastic unravellings.  
In the independent dephasing model considered here, photodetection yields measurement records containing only shot noise, while homodyne and heterodyne detection produce continuous currents that encode (with different signal-to-noise ratios) information about the $\sigma_z$ expectation values of the two qubits.  
These differences will become crucial when we move from independent to engineered, collective collapse operators and use continuous monitoring as a resource to generate and stabilize entangled states.
